\section*{Chapter \ref{ch:b2:overlapping-generations} --- Overlapping Generations}

\begin{solution}{exo:b2:overlapping-generations:1}
% figdata: oadr_summary
Japan, from 8.9 to 51.0 per 100; Nigeria's ratio stayed near 5 (6.0 to 5.4). Japan's births fell and its people live longer, so each cohort reaching
working age is small relative to those reaching 65; Nigeria's births stayed high, so the young keep outnumbering the old.
\end{solution}

\begin{solution}{exo:b2:overlapping-generations:2}
$0.40 \times 0.5 = 20$ per cent of wages.
\end{solution}

\begin{solution}{exo:b2:overlapping-generations:3}
% figdata: eff_summary
Above the line, capital income exceeds gross investment: capital yields more resources than it absorbs, the sign of
efficiency. 128 of 129 economies lie above it, 99 per cent.
\end{solution}

\begin{solution}{exo:b2:overlapping-generations:4}
Gross investment that exceeds the income capital generates, year after year: the economy would be putting more into
capital than it takes out of it.
\end{solution}

\begin{solution}{exo:b2:overlapping-generations:5}
% code: exercises/solutions.py steady_k, ratio_r_to_growth
$k^* = \left(\frac{1 \times (2/3)}{2 \times 1}\right)^{3/2} = (1/3)^{3/2} = 0.192$. Then $R = \tfrac13 k^{-2/3} = \tfrac13 \times (1/3)^{-1} = 1 = 1 + n$: the golden rule.
\end{solution}

\begin{solution}{exo:b2:overlapping-generations:6}
% code: exercises/solutions.py ratio_r_to_growth
$R^*/(1 + n) = \alpha(1 + \beta)/(\beta(1 - \alpha))$: with $\beta = 1$, $0.6/0.7 = 0.857 < 1$, inefficient; with $\beta = 0.6$, $0.48/0.42 = 1.143 > 1$, efficient. In the inefficient case a small pay-as-you-go system lowers capital
toward the golden rule: the first old generation gains its pension, and every later generation gains too, since its
contributions earn $1 + n > R$.
\end{solution}

\begin{solution}{exo:b2:overlapping-generations:7}
% code: exercises/solutions.py contributions_2025
Japan $0.510 \times 0.5 = 25.50$ per cent, Italy 19.85, Korea 14.65, United States 14.25, Nigeria 2.70. The calculation uses people, not contributors and pensioners: it
leaves out who works, who retires before or after 65, pensions financed by funds or taxes, and other spending on the old
such as health care.
\end{solution}

\begin{solution}{exo:b2:overlapping-generations:8}
% code: exercises/solutions.py per_year
$1.01 \times 1.015 - 1 = 2.52$ per cent a year (over thirty years, $(1.01^{30} \times 1.015^{30})^{1/30} - 1$, the same), below the 4 per cent on capital.
\end{solution}

\begin{solution}{exo:b2:overlapping-generations:9}
With $k_{t+1} = Bk_t^\alpha$, the slope is $\alpha Bk^{\alpha - 1} = \alpha k_{t+1}/k_t$, which at the steady state is $\alpha$. A deviation from $k^*$ is multiplied by about $\alpha$ each generation: with $\alpha = 0.3$, 70 per cent
of a gap closes in one generation.
\end{solution}

\begin{solution}{exo:b2:overlapping-generations:10}
The young maximise $\ln(w_t - s_t) + \beta\ln(R_{t+1}s_t)$: the first-order condition $1/(w_t - s_t) = \beta/s_t$ gives $s_t = \beta w_t/(1 + \beta)$, whatever $R_{t+1}$. Firms pay $w_t = (1 - \alpha)Ak_t^\alpha$, the marginal product of labour.
Capital at $t + 1$ is the saving of the $L_t$ young: $K_{t+1} = L_ts_t$, so $k_{t+1} = K_{t+1}/L_{t+1} = s_t/(1 + n) = \frac{\beta(1 - \alpha)A}{(1 + \beta)(1 + n)}k_t^\alpha$.
\end{solution}

\begin{solution}{exo:b2:overlapping-generations:11}
At the steady state $k^{*1-\alpha} = \beta(1 - \alpha)A/((1 + \beta)(1 + n))$, so $R^* = \alpha Ak^{*\alpha - 1} = \alpha(1 + \beta)(1 + n)/(\beta(1 - \alpha))$ and $R^*/(1 + n) = \alpha(1 + \beta)/(\beta(1 - \alpha))$. Inefficiency, $R^* < 1 + n$, holds when
$\alpha(1 + \beta) < \beta(1 - \alpha)$, that is when $\alpha < \beta/(1 + 2\beta)$.
\end{solution}

\begin{solution}{exo:b2:overlapping-generations:12}
The household cares about total resources for old age, $R(s_t + d)$ with its own saving $s_t$. It chooses the same total $s_t + d$ as before, so its own saving falls by $d$
(possible if $d$ does not exceed what it would save anyway). The fund invests $d$ in capital, so capital, $(s_t + d)/(1 + n)$ per young person next period, is unchanged.
\end{solution}

\begin{solution}{pb:b2:overlapping-generations:1}
% code: exercises/solutions.py brenland; listings/olg.py
\textbf{Part I.}
1. $w = 0.7 \times 10 \times k^{0.3} = 7k^{0.3}$.
2. $R = 0.3 \times 10 \times k^{-0.7} = 3k^{-0.7}$.
3. $\beta/(1 + \beta) = 0.6/1.6 = 0.375$.
4. $k_{t+1} = 0.375 \times 7k_t^{0.3}/1.2 = 2.1875k_t^{0.3}$.
5. Next period's capital is shared among the next generation of young, $1.2$ times as many.
\textbf{Part II.}
6. $k^* = 2.1875^{1/0.7} = 3.059$.
7. $w = 7 \times 3.059^{0.3} = 9.790$; $R = 3 \times 3.059^{-0.7} = 1.371$.
8. $1.371^{1/30} - 1 = 1.06$ per cent a year.
9. Saving $0.375 \times 9.790 = 3.671$; consumption when young $9.790 - 3.671 = 6.119$; when old $1.371 \times 3.671 = 5.035$.
10. $R = 1.2$ at $k = (3/1.2)^{1/0.7} = 3.702$. Since $k^* = 3.059 < 3.702$ and $R^* > 1.2$, Brenland is efficient.
11. $1.06$ per cent against $1.2^{1/30} - 1 = 0.61$ per cent a year.
\textbf{Part III.}
12. $(1 + n)d = 1.2$.
13. $1.2$ a generation, 0.61 per cent a year.
14. $k^* = 2.059$.
15. $w = 7 \times 2.059^{0.3} = 8.694$; $R = 3 \times 2.059^{-0.7} = 1.810$.
16. The young save less for their old age, since the pension provides part of it, and capital is what the young save:
$s = 1.2 \times 2.059 = 2.471$ against 3.671.
17. No: the fund invests the contributions in capital, and the young cut their own saving by the same amount.
\textbf{Part IV.}
18. A pension of 1.2 thousand dalens each, 1,200 dalens, for nothing.
19. The pension's value at the time of the contribution is $1.2/1.810 = 0.663$; she loses $1 - 0.663 = 0.337$ per contribution of 1.
20. Ending it means either that the young stop paying, so the old of that period lose the pension they paid for, or that
someone pays both for the current old and for their own old age.
21. $0.30 \times 0.5 = 15$ per cent.
22. $0.45 \times 0.5 = 22.5$ per cent.
23. $0.15/0.45 = 0.333$, one third.
24. $0.35 \times 0.5 = 17.5$ per cent.
25. Without the system, capital per young person is 3.06 thousand dalens; with a contribution of 1 thousand, 2.06, a third lower. The
first old generation receives a windfall of 1,200 dalens each, which later generations pay for with a contribution that
returns 1.2 against 1.81 on capital.
\end{solution}
